% ===================================================================
%  সারণি নং 11 — নগর ও তার ইমারত। — আগুন।
%
%  Ceci n'est PAS une transcription : c'est une traduction, faite en
%  2026, en bengali d'aujourd'hui. Voir texto/bn/10-sarani-01.tex
%  pour la consigne, qui vaut pour les seize fichiers.
% ===================================================================

%%K t11-apar-1 apar
\VUsaut{2.0mm}
\VUcentre{13.2pt}{90}{তৃতীয় পর্যায়}
\VUsaut{3.0mm}
\VUfilet{16mm}
\VUsaut{5.6mm}
\VUtitre{15.0pt}{60}{সারণি নং 11}
\VUsaut{1.4mm}
\VUpk{11.4pt}{40}{নগর ও তার ইমারত। --- আগুন।}
\VUsaut{2.2mm}

%%K t11-tit-1 sub
\VUpk{10.2pt}{40}{বাইরের পাড়া।}
\VUsaut{3.0mm}

%%K t11-c1-01-1 p
1. --- আমি মহাসাগর থেকে একশো কিলোমিটার দূরে একটি বড় নগরে থাকি, যা তবু
প্রথম শ্রেণির \VUgras{বন্দর} \textsuperscript{(1)}। যে নদী তার পাশ দিয়ে বয় তা
চারশো মিটার চওড়া এবং একটি খুব সুন্দর নোঙরগাহ গড়ে।

%%K t11-c1-02-1 p
2. --- \VUgras{ঘাটের} \textsuperscript{(2)} ধারের সমস্ত অংশ একেবারে সামুদ্রিক ও
শিল্পের চেহারা রাখে, আর এমন একটি \VUgras{বাইরের পাড়া} \textsuperscript{(3)}
গড়ে যেখানে কেবল নাবিক ও মজুরেরা থাকে। এরা \VUgras{কারখানায়}
\textsuperscript{(4)} ও \VUgras{গ্যাসঘরে} \textsuperscript{(5)} কাজ করে, যার উঁচু
\VUgras{চিমনি} \textsuperscript{(6)} ও \VUgras{গ্যাস-ট্যাংক} দূর থেকেই চোখে পড়ে।

%%K t11-c1-03-1 p
3. --- মধ্যযুগে নগর দুর্গপ্রাচীরে ঘেরা ছিল, আর তার প্রাচীরের কিছু অবশেষ এখনও
মেলে, বিশেষ করে সেই \VUgras{ফটক} \textsuperscript{(8)} যা পুরনো \VUgras{গড়-দুর্গের}
\textsuperscript{(7)}, আর যার দুই পাশে তার দুটি \VUgras{মিনার} \textsuperscript{(9)},
যা এখন \VUgras{কারাগারের} কাজে লাগে।

%%K t11-c1-04-1 p
4. --- এই সমস্ত \VUgras{পাড়ায়}, যা খুব পুরনো লাগে, কেবল একটি ইমারত
তুলনামূলকভাবে নতুন: \VUgras{শুল্কঘর} \textsuperscript{(10)}। সেটি ঘাট ও সেই ছোট
\VUgras{গলির} \textsuperscript{(11)} মাঝে দাঁড়ানো যা পুরনো \VUgras{নগরভবন}
\textsuperscript{(12)} পর্যন্ত যায়। সেই শেষ ইমারতটিকে সহজেই চেনা যায় সেই বিশাল
\VUgras{ঘণ্টা-মিনারে} \textsuperscript{(13)} যা পুরনো নগরের উপরে ছেয়ে থাকে।

%%K t11-c1-05-1 p
5. --- গলির অন্য দিকে শুল্কঘরেরই ধাঁচের একটি ইমারত \VUgras{দাঁড়িয়ে}: সেটি
\VUgras{বাণিজ্যভবন} \textsuperscript{(14)}। তার একটি মুখ একটি ছোট চত্বরের দিকে
খোলে যেখানে \VUgras{প্রশাসনভবন} \textsuperscript{(15)}। আমাদের নগর, বস্তুত,
রাজধানী না হয়েও একটি গুরুত্বপূর্ণ জেলাসদর।

%%K t11-tit-2 sub
\VUsaut{5.0mm}
\VUpk{10.2pt}{40}{নগরের উপরের অংশ।}
\VUsaut{3.4mm}

%%K t11-c1-06-1 p
6. --- ছাউনি, যা বেশ বড়, বিশাল ও খুব স্বাস্থ্যকর \VUgras{ব্যারাকে}
\textsuperscript{(16)} থাকে; সেগুলি \VUgras{নগর-উদ্যানের} \textsuperscript{(17)}
জাফরি পর্যন্ত ছড়ানো। যে \VUgras{চওড়া রাস্তা} \textsuperscript{(19)} এই
উদ্যান পর্যন্ত যায়, তা \VUgras{সঞ্চয়-ব্যাংকের} \textsuperscript{(18)} পিছন দিয়ে
গিয়ে \VUgras{মহাগির্জার} \textsuperscript{(20)} চত্বরে বেরোয়।

%%K t11-c1-07-1 p
7. --- এই সব ইমারত একটি ছোট পাহাড়ে সিঁড়ির মতো উঠেছে, যেখানে আমাদের বড়
\VUgras{মাধ্যমিক বিদ্যালয়ও} \textsuperscript{(21)} আছে।

%%K t11-c1-07-2 p
আপনি যদি সেই হালকা ঢালু রাস্তা দিয়ে নামেন যা বড় রাস্তার সঙ্গে মেলে, তবে আপনি
\VUgras{আদালত} \textsuperscript{(22)} পর্যন্ত পৌঁছান, যার সামনে একটি মনোরম
\VUgras{সর্বজনীন বাগান} \textsuperscript{(26)}। এই \VUgras{চত্বর} খুব বড় নয়,
কিন্তু নগরের মাঝখানে সেটি সবুজ ও ঠান্ডার মরূদ্যান গড়ে। তাই নগরবাসীরা তাতে
খুবই আসে, যাদের \VUgras{কফিঘরের} \textsuperscript{(25)} ছাউনির নিচে বসে সেই
সামরিক ব্যান্ড শুনতে ভালো লাগে যা বৃহস্পতিবার ও রবিবার লন ও কেয়ারির মাঝে গড়া
\VUgras{বাদ্যমঞ্চে} \textsuperscript{(27)} বাজে।

%%K t11-c1-08-1 p
8. --- সেই লনগুলিরই একটিতে ব্রোঞ্জের একটি \VUgras{মূর্তি} \textsuperscript{(28)}
দাঁড়ানো, আমাদের এক নগরবাসীর স্মৃতিতে তোলা স্মারক, যিনি নিজের জন্মভূমির বড়
সেবা করেছিলেন।

%%K t11-tit-3 sub
\VUsaut{4.0mm}
\VUpk{10.2pt}{40}{যাতায়াত।}
\VUsaut{3.4mm}

%%K t11-c1-09-1 p
9. --- আমাদের নগর এখনও বিদ্যুতে আলোকিত নয়; তাতে কেবল \VUgras{গ্যাসের বাতি}
\textsuperscript{(29)}। কিন্তু \VUgras{স্থানীয় কাগজ} এই নিয়ে অভিযান চালাচ্ছে
যে আলোর সেই ব্যবস্থা শোধরানো হোক, যেমন ট্রামের \VUgras{টানার ব্যবস্থা} আগেই
শোধরানো হয়েছে। \VUgras{ঘোড়ায়} \VUgras{টানা} পুরনো গাড়ির জায়গায় এখন
কাজের \VUgras{বৈদ্যুতিক ট্রাম} \textsuperscript{(31)} এসে গেছে, যা খুব
\VUgras{কম} ভাড়ায় যাত্রীদের \VUgras{নগরের} এক প্রান্ত থেকে \VUgras{অন্য
প্রান্তে} তাড়াতাড়ি পৌঁছে দেয়।

%%K t11-c1-10-1 p
10. --- আদালতের পাশে \VUgras{ব্যাংক} \textsuperscript{(32)}, যেখানে বাণিজ্যিক
হুন্ডি ভাঙানো হয়। \VUgras{ব্যাংকের বার্তাবাহকেরা} \textsuperscript{(33)} লোকের
বাড়ি বাড়ি গিয়ে বকেয়া আদায় করে।

%%K t11-tit-4 sub
\VUsaut{4.0mm}
\VUpk{10.2pt}{40}{শিল্প ও বিদ্যা।}
\VUsaut{3.4mm}

%%K t11-c1-11-1 p
11. --- পাশের সেই সুন্দর ইমারতটি \VUgras{সংগ্রহশালা}। তাতে একসঙ্গে নগরের
\VUgras{গ্রন্থাগার} \textsuperscript{(34)}, কয়েকটি সুন্দর \VUgras{প্রাকৃতিক
ইতিহাসের সংগ্রহ} \textsuperscript{(35)} (একটি প্রাকৃতিক ইতিহাস সংগ্রহশালা) ও
একটি সুঠাম \VUgras{চিত্রশালা} \textsuperscript{(36)}, যা একটি চমৎকার স্থায়ী
\VUgras{প্রদর্শনী} গড়ে।

%%K t11-c1-11-2 p
সেটি সপ্তাহে দুবার জনসাধারণের জন্য খোলে, কিন্তু দর্শকেরা যে কোনো দিন সেটি
দেখতে পারেন, \VUgras{রক্ষককে} \textsuperscript{(37)} একটু বকশিশ দিয়ে।
\VUgras{চিত্রকরেরা} \textsuperscript{(38)} সেখানে কাজ করতে আসেন। তাঁদের নিজেদের
তেপায়ার \VUgras{সামনে} দেখা যায়, হাতে \VUgras{রঙের পাত} নিয়ে, বিখ্যাত ছবির
নকল করতে।

%%K t11-c1-12-1 p
12. --- সংগ্রহশালার পিছনের উঁচু জায়গায় সেই \VUgras{গম্বুজ}
\textsuperscript{(40)} যা \VUgras{হাসপাতালের} \textsuperscript{(39)}, আর সেই
\VUgras{শিক্ষক-প্রশিক্ষণ বিদ্যালয়ের} \textsuperscript{(41)} ইমারত কষ্টে চোখে
পড়ে যেখানে আমাদের নগর-বিদ্যালয়ের শিক্ষকেরা প্রশিক্ষিত হয়েছিলেন। রঙ্গমঞ্চের
চত্বরে একটি \VUgras{ডাকঘর} \textsuperscript{(43)}, যা একটি \VUgras{ব্যক্তিগত
বাড়িতে} \textsuperscript{(42)} বসানো।

%%K t11-tit-5 sub
\VUsaut{5.0mm}
\VUpk{11.4pt}{40}{আগুন।}
\VUsaut{3.0mm}

%%K t11-tit-6 sub
\VUpk{10.2pt}{40}{আগুনের হাঁক।}
\VUsaut{3.4mm}

%%K t11-c2-01-1 p
1. --- নগরে ভয়ংকর খবর ছড়াচ্ছে: রঙ্গমঞ্চে এইমাত্র আগুন লেগেছে! আমি যখন
\VUgras{চত্বরে} \textsuperscript{(96)} পৌঁছাই, ভিড় \VUgras{ফুটপাতে}
\textsuperscript{(46)} ও \VUgras{রাস্তায়} \textsuperscript{(47)} জমে গেছে।
\VUgras{ডাকের কেরানিরা} \textsuperscript{(44)} ও \VUgras{ডাকপিয়নেরা}
\textsuperscript{(45)} ডাকঘর থেকে বেরিয়ে আসছে। একজন \VUgras{ট্রামচালককে}
\textsuperscript{(48)} নিজের ট্রাম থামাতে হয়েছে যাতে নগরের একটি
\VUgras{রোগীবাহী গাড়ি} \textsuperscript{(50)} বেরোতে পারে, যা একজন
\VUgras{শল্যচিকিৎসক} \textsuperscript{(52)} ও একজন \VUgras{সেবিকা}
\textsuperscript{(51)} নিয়ে আসছে, একজন \VUgras{আহতের} \textsuperscript{(53)}
দেখাশোনার জন্য, যাকে \VUgras{স্ট্রেচারে} \textsuperscript{(54)}
\VUgras{খবরের কাগজের খুপরির} \textsuperscript{(55)} কাছে আনা হয়েছে।

%%K t11-c2-02-1 p
2. --- কসরত থেকে ফেরা পদাতিক সেনার একটি দল থেমে গেছে, যাতে ঘোড়ায় সওয়ার
\VUgras{নগররক্ষীদের} \textsuperscript{(61)} সঙ্গে শৃঙ্খলা রাখতে সাহায্য করে।
\VUgras{ঘোড়সওয়ারদের} ঘোড়া দু পায়ে দাঁড়িয়ে পড়ে; তারা \VUgras{সৈনিক}
\textsuperscript{(57)} বা \VUgras{জঁদার্মদের} \textsuperscript{(63)} চেয়ে ভালো
\VUgras{দর্শকদের} \textsuperscript{(56)} পিছনে ঠেলতে পারে। জঁদার্মেরা এমনিতেও খুব
ব্যস্ত। তাদের একজন \VUgras{পুলিশের সিপাহিদের} \textsuperscript{(64)} বলছে যে আর
একজন \VUgras{আহতকে} \textsuperscript{(65)} কোথায় নিয়ে যেতে হবে, এক বীর দমকলকর্মী
যে সিঁড়ি থেকে পড়ে গেছে।

%%K t11-c2-03-1 p
3. --- \VUgras{অফিসার} \textsuperscript{(66)} ঠিক ঠিক বলছেন যে যে \VUgras{বাষ্প-পাম্প}
\textsuperscript{(67)} এইমাত্র আসছে তা কোথায় রাখতে হবে। সেই শক্তিশালী যন্ত্রের
অপেক্ষায় তিনি একটি \VUgras{হাতপাম্প} \textsuperscript{(68)} বসিয়ে দিয়েছেন,
যার লম্বা \VUgras{নল} \textsuperscript{(69)} আগুনের উপরে জলের ধারা ঢালছে। ছয়জন
দমকলকর্মী সেটি চালাচ্ছে, যতক্ষণ তাদের সঙ্গী \VUgras{টোঁটি}
\textsuperscript{(70)} ধরে \VUgras{ধারাকে} \textsuperscript{(71)} আগুনের মাঝে
তাক করছে।

%%K t11-c2-04-1 p
4. --- রঙ্গমঞ্চের অন্য দিকে বড় \VUgras{মই} \textsuperscript{(72)} খাড়া করা
হয়েছে। তা দমকলকর্মীদের সেই শিখার কাছে পৌঁছাতে দেয় যা কালো \VUgras{ধোঁয়ার}
\textsuperscript{(75)} ঘূর্ণির মাঝে উঠছে।

%%K t11-tit-7 sub
\VUsaut{5.0mm}
\VUpk{10.2pt}{40}{বীরত্বের কাজ।}
\VUsaut{3.4mm}

%%K t11-c2-05-1 p
5. --- \VUgras{আগুন} \textsuperscript{(74)} \VUgras{রঙ্গমঞ্চের}
\textsuperscript{(73)} প্রবেশকক্ষে লেগেছিল, সেই \VUgras{অপেরার} মহড়ার সময় যার
প্রথম \VUgras{প্রদর্শনী} কাল হওয়ার কথা ছিল। সৌভাগ্যক্রমে সভাকক্ষে কেবল কয়েকজন
\VUgras{দর্শক} \textsuperscript{(76)} ছিলেন। বেচারারা \VUgras{বারান্দায়}
\textsuperscript{(77)} আশ্রয় নিয়েছেন, যেখানে বীর \VUgras{দমকলকর্মীরা}
\textsuperscript{(86)} মই ও দড়ি নিয়ে তাঁদের সাহায্যে আসছে।

%%K t11-c2-06-1 p
6. --- \VUgras{তোরণের} \textsuperscript{(78)} নিচে, যেখানে \VUgras{নাট্যবিজ্ঞপ্তি}
\textsuperscript{(79)} প্রদর্শনীর ঘোষণা করে, \VUgras{বাদকদের}
\textsuperscript{(81)} পালাতে দেখা যাচ্ছে, যাঁরা \VUgras{বাদ্যদলের}
\textsuperscript{(80)}, নিজেদের বাদ্যযন্ত্র তুলে: \VUgras{বেহালা}
\textsuperscript{(81)}, \VUgras{বাঁশি} \textsuperscript{(82)}, \VUgras{চেলো}
\textsuperscript{(83)}, \VUgras{ট্রম্বোন} \textsuperscript{(84)}, \VUgras{ঢোল}
\textsuperscript{(85)}, ইত্যাদি। \VUgras{আসবাবের} \textsuperscript{(87)} একটি অংশ
বাঁচানোর চেষ্টা হচ্ছে।

%%K t11-c2-07-1 p
7. --- \VUgras{অভিনেতা} \textsuperscript{(89)} ও \VUgras{অভিনেত্রীদের}
\textsuperscript{(88)}, গায়ক ও গায়িকাদের, নিজেদের মঞ্চের \VUgras{পোশাকেই}
\textsuperscript{(90)} পালাতে হয়েছে। মুখ্য গায়ক একজন \VUgras{সাংবাদিককে}
\textsuperscript{(92)} বোঝাচ্ছেন যে এটি কীভাবে ঘটল। \VUgras{সংবাদদাতা} প্রথম
মুহূর্ত থেকেই কর্তৃপক্ষের সঙ্গে ছুটে এসেছিলেন: \VUgras{প্রশাসক}
\textsuperscript{(91)}, \VUgras{নগরপাল} \textsuperscript{(93)}, \VUgras{পুলিশ
অধীক্ষক} \textsuperscript{(94)} ও একজন \VUgras{বিচারক} \textsuperscript{(95)},
যাঁরা তদন্ত করবেন কার দায়িত্ব।
\VUsaut{6.0mm}
\VUfilet{22mm}
